\documentclass[12pt]{article}
\usepackage[utf8]{inputenc}
\usepackage[T1]{fontenc}
\usepackage{hyperref}
\title{Tinigua Language Resources for K-12 Curriculum Development}
\author{Compiled by Claude Sonnet 5 (claude-sonnet-5) for Andr\'as Kornai}
\date{}
\begin{document}
\maketitle
\section{Introduction}
Tinigua (ISO 639-3 tit; Glottocode tini1245) is a critically endangered language of Meta Department, Colombia, historically spoken between the upper Guayabero and Yari rivers. It belonged to a small family with the extinct Pamigua. Search results indicate that the language is effectively down to a single known speaker (Sixto Mu\~noz), who reportedly did not pass it on to his children; this is lower than the figure of about 2 in the reference description. No school programme, primer or native-language curriculum material was found.
\section{Primary Educational Resources}
None found. An educational video/resource page on Tinigua exists in the Argentine ``Guardianes de la lengua'' series, but its content was not examined:
\url{https://www.educ.ar/recursos/157164/guardianes-de-la-lengua-tinigua}
\section{Linguistic and Reference Materials}
General descriptions of the language and family: \url{https://en.wikipedia.org/wiki/Tinigua_language} and \url{https://en.wikipedia.org/wiki/Tiniguan_languages}. A numerals data sheet is at \url{https://lingweb.eva.mpg.de/channumerals/Tinigua.htm}. The main early source mentioned is Castellv\'i (1940), \emph{La lengua tinigua}; no URL was seen for it.
\section{Religious and Community Resources}
None found.
\section{Digital and Audio Resources}
Language-status pages: \url{https://ar.wal.unesco.org/en/countries/colombia/languages/tinigua} and \url{https://www.endangeredlanguages.com/lang/2146?hl=zh-cn}. Press context: \url{https://colombiareports.com/colombias-indigenous-languages-disappearing/}. No audio corpus was found.
\section{Curriculum Development Context}
With at most one living speaker, a K-12 curriculum in Tinigua is not currently feasible; any work would be documentation and revitalization from archival sources, in consultation with the community and the remaining speaker.
\end{document}